
\chapter{Trabalhos Relacionados}
\label{cap:relacionados}

Este capítulo revisa os trabalhos relacionados em quatro frentes: a aplicação de modelos de linguagem de grande escala à engenharia de requisitos, o reconhecimento de fala para a transcrição de reuniões, a avaliação da qualidade de histórias de usuário e os sistemas multiagente voltados à engenharia de software. Ao final, sintetizam-se as lacunas identificadas na literatura que motivam esta pesquisa.

\section{Modelos de Linguagem de Grande Escala na Engenharia de Requisitos}

A aplicação de modelos de linguagem de grande escala à engenharia de requisitos tem atraído interesse crescente \cite{hort2025survey}. Motger et al.~\cite{motger2025share} conduziram um estudo de mapeamento sistemático que identificou 62 conjuntos de dados publicamente disponíveis para pesquisa em ER baseada em LLMs, constatando que apenas 2 deles tratam de tarefas de elicitação e que 58 dos 62 são exclusivamente em inglês, uma escassez que motiva diretamente a validação em língua portuguesa apresentada neste trabalho.

No que diz respeito à geração de artefatos, Delgado et al.~\cite{delgado2024aidriven} compararam heurísticas baseadas em \textit{N-gram} com o GPT-3 para a geração automática de histórias de usuário a partir de especificações textuais já existentes, reportando ROUGE-N\,=\,0,46 e BERTScore\,=\,0,69 para a abordagem baseada em LLM. Ren et al.~\cite{ren2024combining} demonstraram, ainda, que a técnica de \textit{Chain-of-Thought} permite que LLMs extraiam diagramas de classes conceituais a partir de histórias de usuário. Naimi et al.~\cite{naimi2024automating} utilizaram LLMs para gerar a descrição textual de casos de uso a partir de diagramas UML codificados como XML, reduzindo o tempo de documentação e melhorando a consistência. Em conjunto, esses trabalhos demonstram a viabilidade dos LLMs para a geração de artefatos de ER, mas compartilham um pré-requisito comum: a entrada textual estruturada. Nenhum deles trata da extração de requisitos a partir do áudio não estruturado de reuniões.

\section{Reconhecimento de Fala para Transcrição de Reuniões}

O reconhecimento automático de fala (do inglês \textit{Automatic Speech Recognition}, ASR) é a tecnologia habilitadora da elicitação de requisitos baseada em áudio. Radford et al.~\cite{whisper2022} introduziram o Whisper, um modelo treinado por supervisão fraca em larga escala que demonstra transcrição multilíngue robusta sob diversas condições acústicas, incluindo falantes não nativos e ruído de fundo, características comuns em reuniões reais de partes interessadas. Kasakowskij e Haake~\cite{kasakowskij2025talk2text} desenvolveram o T\textsuperscript{3}~Talk2Text, um sistema de ASR em tempo quase real para reuniões virtuais em grupo, que captura transcrições de todo o grupo para a análise posterior de padrões de comunicação e colaboração.

Apesar desses avanços, o ASR em reuniões de elicitação de requisitos enfrenta desafios específicos do domínio: vocabulário técnico, fala sobreposta e siglas específicas de projeto que modelos de propósito geral podem transcrever incorretamente. Tais erros não permanecem isolados: eles se propagam para os agentes a jusante, um modo de falha que se analisa empiricamente na QP5.

\section{Qualidade de Histórias de Usuário e o Framework INVEST}

As histórias de usuário são o artefato de requisitos predominante no desenvolvimento ágil, seguindo o formato \textit{``Como [papel], quero [objetivo] para [benefício]''}~\cite{lucassen2016qus}. Wake~\cite{wake2003invest} propôs as heurísticas INVEST (\textit{Independent}, \textit{Negotiable}, \textit{Valuable}, \textit{Estimable}, \textit{Small} e \textit{Testable}) como seis critérios para avaliar se uma história de usuário é acionável e bem delimitada.

Lucassen et al.~\cite{lucassen2016qus} formalizaram o \textit{framework} de \textit{Quality User Story} (QUS), composto por 13 critérios que cobrem sintaxe, semântica e pragmática, e introduziram a AQUSA, uma ferramenta baseada em processamento de linguagem natural que detecta defeitos de qualidade. Avaliada sobre 1.023 histórias de 18 organizações de software, a AQUSA demonstra a viabilidade da avaliação automatizada de qualidade, mas requer histórias escritas manualmente como entrada. Xu et al.~\cite{xu2023requirement} combinaram o INVEST com um modelo de sete critérios para automatizar a avaliação de qualidade em \textit{pipelines} ágeis. Este trabalho estende essa linha ao gerar histórias de usuário a partir do áudio de reuniões e ao avaliá-las frente ao INVEST por meio de um protocolo de LLM como juiz validado contra avaliadores humanos.

\section{Sistemas Multiagente para Engenharia de Software}

Arquiteturas multiagente emergiram como um paradigma produtivo para a automação de tarefas complexas de engenharia de software. O ChatDev~\cite{qian2024chatdev} introduziu um \textit{framework} baseado em diálogo no qual agentes LLM especializados colaboram por meio de conversas estruturadas ao longo das fases de projeto, codificação e teste, demonstrando que a comunicação linguística pode unificar agentes autônomos. O MetaGPT~\cite{hong2024metagpt} codificou Procedimentos Operacionais Padronizados (SOPs) em sequências de \textit{prompts}, atribuindo papéis de engenharia de software (gerente de produto, arquiteto, engenheiro e responsável por qualidade) a agentes distintos, alcançando resultados do estado da arte em \textit{benchmarks} de geração de código. Luo et al.~\cite{luo2024repoagent} aplicaram um \textit{pipeline} agêntico à documentação, com três agentes especializados produzindo e mantendo a documentação de código em nível de repositório, superando as linhas de base produzidas por humanos. Há, ainda, esforços rumo a plataformas unificadas de desenvolvimento multiagente~\cite{sami2024experimentingmultiagent} e à automação de fluxos ágeis completos, da criação de histórias de usuário à geração e revisão de código~\cite{sanwal2024autonomous}. Um levantamento recente de Cai et al.~\cite{cai2025designing} analisou 94 artigos sobre sistemas multiagente baseados em LLMs para tarefas de engenharia de software, constatando que a \textit{Cooperação Baseada em Papéis} (\textit{Role-Based Cooperation}) é o padrão de projeto mais comumente adotado, o mesmo padrão que fundamenta o \textit{pipeline} proposto neste trabalho.

O trabalho mais próximo do nosso é o MARE~\cite{jin2024mare}, um \textit{framework} de colaboração multiagente para todo o processo de ER. O MARE decompõe a ER em quatro tarefas (elicitação, modelagem, verificação e especificação) distribuídas entre cinco agentes especializados, superando as linhas de base do estado da arte em até 15,4\% em F1. Contudo, o MARE opera a partir de uma \textit{ideia inicial} (\textit{rough idea}) textual fornecida por um humano e simula a elicitação por meio de diálogo entre agentes; ele não processa o áudio de reuniões e, portanto, não pode ser implantado em um cenário pós-reunião totalmente automatizado. O sistema proposto neste trabalho aborda precisamente essa lacuna.

\section{Lacunas Identificadas}

A literatura revisada revela quatro lacunas que motivam este trabalho:

\begin{enumerate}
    \item[1.] \textbf{Ausência de \textit{pipelines} de áudio para SRS.} As ferramentas comerciais existentes \cite{fellow2025,jamie2025} produzem anotações de reunião de propósito geral, e não artefatos formais de ER. Os sistemas de ER baseados em LLMs, como o MARE \cite{jin2024mare}, processam texto estruturado, e não o áudio conversacional bruto.

    \item[2.] \textbf{Foco em subtarefas isoladas.} O trabalho anterior com LLMs em ER trata de tarefas individuais (geração de histórias de usuário \cite{delgado2024aidriven}, descrição de casos de uso \cite{naimi2024automating}, modelagem conceitual \cite{ren2024combining}) sem integrar a transcrição, a identificação de personas e a geração de documentos em um único \textit{pipeline}.

    \item[3.] \textbf{Falta de validação no mundo real.} A maioria dos estudos utiliza conjuntos de dados sintéticos ou \textit{corpora} controlados; poucos avaliam o resultado frente a um documento de requisitos oficial produzido manualmente em um projeto de software real.

    \item[4.] \textbf{Propagação de erro não analisada.} Em \textit{pipelines} agênticos, os erros de transcrição se propagam em cascata para os artefatos a jusante. Esse modo de falha é reconhecido, mas não é medido empiricamente em nenhum dos trabalhos relacionados.
\end{enumerate}

Este trabalho aborda as quatro lacunas por meio de um \textit{pipeline} de quatro agentes que transforma o áudio de reuniões em documentação de requisitos estruturada, validado sobre um projeto de software governamental real e com a medição quantitativa da propagação de erro entre agentes.
